\documentclass[11pt,a4paper]{article}
\usepackage[T1]{fontenc}
\usepackage[utf8]{inputenc}
\usepackage{amsmath,amssymb,amsthm}
\usepackage[margin=1in]{geometry}
\usepackage[colorlinks=true,linkcolor=blue,urlcolor=blue]{hyperref}
\theoremstyle{plain}
\newtheorem{theorem}{Theorem}
\newtheorem{lemma}[theorem]{Lemma}
\newtheorem{proposition}[theorem]{Proposition}
\newtheorem{corollary}[theorem]{Corollary}
\theoremstyle{definition}
\newtheorem*{remark}{Remark}

\title{A closed form for the $p$-adic valuation of OEIS A129365,\\
and a proof that its terms are integers}
\author{Adrian Perez Fontelles\\ \small Independent researcher}
\date{23 August 2026}
\begin{document}
\maketitle

\begin{abstract}
OEIS A129365 is the quotient
$a(n)=\bigl(\prod_{j,k\le n}\gcd(j,k)\bigr)\big/\bigl(\prod_{k=1}^{n}(\lfloor
n/k\rfloor!)^{k}\bigr)$. In April 2007 P.~Bala listed four conjectures about it, the first
being simply that $a(n)$ is always an integer. We prove the exact closed form
\[
  \operatorname{ord}_{p}a(n)=\sum_{i\ge1}B\bigl(\lfloor n/p^{i}\rfloor\bigr),
  \qquad B(M)=\sum_{k=1}^{M}\bigl(M\bmod k\bigr)\ \ (\text{A004125}),
\]
for every prime $p$. Since a sum of remainders is nonnegative, every exponent is
nonnegative and Conjecture A follows at once. The proof is three steps: the known
factorisation of the numerator, the nested-floor identity
$\lfloor n/(kp^{i})\rfloor=\lfloor\lfloor n/p^{i}\rfloor/k\rfloor$ for the denominator, and
the observation that $M^{2}-\sum_{k\le M}k\lfloor M/k\rfloor$ is exactly
$\sum_{k\le M}(M\bmod k)$ because $M^{2}=\sum_{k\le M}M$.
\end{abstract}

\section{The sequence and the conjecture}

OEIS A129365 (Peter Bala, 13 April 2007) has offset 1 and is defined by
$a(n)=\mathrm{A092287}(n)/\mathrm{A129364}(n)$, equivalently by the formula recorded on the
entry
\begin{equation}\label{eq:def}
  a(n)=\frac{\displaystyle\prod_{j=1}^{n}\prod_{k=1}^{n}\gcd(j,k)}
            {\displaystyle\prod_{k=1}^{n}\bigl(\lfloor n/k\rfloor!\bigr)^{k}},
\end{equation}
with $a(1),\dots,a(10)=1,1,1,1,1,2,2,2,6,48$. The entry carries four conjectures, labelled
A--D. This note proves the first:

\begin{quote}
\textbf{Conjectures:} A) $a(n)$ is always an integer.
\end{quote}

Conjectures B and C were addressed by T.~Adamczewski (arXiv:2608.11941, August 2026), as
annotated on the entry by M.~De Vlieger on 20 August 2026. Conjecture D is the same closed
form evaluated at $n\mapsto np$. As of the ``Last modified'' line on the live entry
(23 August 2026), A and D are still recorded as unproven.

Throughout $p$ is a prime, $v_{p}(N)$ is the exponent of $p$ in $N$ (the entry's
$\operatorname{ord}_{p}$), and
\[
  B(M)=\sum_{k=1}^{M}\bigl(M\bmod k\bigr)
\]
is A004125, the sum of remainders, with $B(0)=0$.

\section{Three ingredients}

\begin{lemma}[numerator]\label{lem:num}
$\displaystyle v_{p}\Bigl(\prod_{j=1}^{n}\prod_{k=1}^{n}\gcd(j,k)\Bigr)
=\sum_{i\ge1}\lfloor n/p^{i}\rfloor^{2}$.
\end{lemma}

\begin{proof}
$v_{p}(\gcd(j,k))=\min(v_{p}(j),v_{p}(k))$ counts the $i\ge1$ with $p^{i}\mid j$ and
$p^{i}\mid k$. Summing over $j,k\le n$ and exchanging the order of summation, the two
divisibility conditions are independent, so the count at level $i$ is
$\bigl(\#\{j\le n:p^{i}\mid j\}\bigr)^{2}=\lfloor n/p^{i}\rfloor^{2}$. (This is the
square-case formula of A092287, confirmed on that entry by C.~R. Greathouse IV in April
2013.)
\end{proof}

\begin{lemma}[nested floors]\label{lem:floor}
For positive integers $n,k,q$, $\lfloor n/(kq)\rfloor=\lfloor\lfloor n/q\rfloor/k\rfloor$.
\end{lemma}

\begin{proof}
Standard: for a real $x$ and a positive integer $k$,
$\lfloor x/k\rfloor=\lfloor\lfloor x\rfloor/k\rfloor$; apply it to $x=n/q$.
\end{proof}

\begin{lemma}[remainder identity]\label{lem:rem}
For every integer $M\ge0$,
$\displaystyle M^{2}-\sum_{k=1}^{M}k\lfloor M/k\rfloor=\sum_{k=1}^{M}\bigl(M\bmod k\bigr)
=B(M)$.
\end{lemma}

\begin{proof}
$M^{2}=\sum_{k=1}^{M}M$, and $M-k\lfloor M/k\rfloor=M\bmod k$ termwise.
\end{proof}

\section{The closed form}

\begin{theorem}\label{thm:main}
For every $n\ge1$ and every prime $p$,
\[
  v_{p}\bigl(a(n)\bigr)=\sum_{i\ge1}B\bigl(\lfloor n/p^{i}\rfloor\bigr),
\]
where $a(n)$ is given by \eqref{eq:def}. All terms with $p^{i}>n$ vanish, so the sum is
finite.
\end{theorem}

\begin{proof}
Write $M_{i}=\lfloor n/p^{i}\rfloor$. By Legendre's formula and Lemma~\ref{lem:floor},
\[
  v_{p}\Bigl(\prod_{k=1}^{n}\bigl(\lfloor n/k\rfloor!\bigr)^{k}\Bigr)
  =\sum_{k=1}^{n}\sum_{i\ge1}k\bigl\lfloor\lfloor n/k\rfloor\big/p^{i}\bigr\rfloor
  =\sum_{i\ge1}\sum_{k=1}^{n}k\bigl\lfloor n/(kp^{i})\bigr\rfloor
  =\sum_{i\ge1}\sum_{k=1}^{M_{i}}k\lfloor M_{i}/k\rfloor,
\]
where in the last step Lemma~\ref{lem:floor} was used again in the form
$\lfloor n/(kp^{i})\rfloor=\lfloor M_{i}/k\rfloor$, and the terms with $k>M_{i}$ were
dropped because $\lfloor M_{i}/k\rfloor=0$ there. Subtracting this from Lemma~\ref{lem:num}
and applying Lemma~\ref{lem:rem} at each level $i$,
\[
  v_{p}\bigl(a(n)\bigr)=\sum_{i\ge1}\Bigl(M_{i}^{2}-\sum_{k=1}^{M_{i}}k\lfloor M_{i}/k
  \rfloor\Bigr)=\sum_{i\ge1}B(M_{i}).\qedhere
\]
\end{proof}

\begin{corollary}[Conjecture A]\label{cor:A}
$a(n)$ is an integer for every $n\ge1$.
\end{corollary}

\begin{proof}
Each $M\bmod k$ is a nonnegative integer, so $B(M)\ge0$ and hence $v_{p}(a(n))\ge0$ for
every prime $p$ by Theorem~\ref{thm:main}. A positive rational with nonnegative valuation
at every prime is an integer.
\end{proof}

\begin{remark}[the other three conjectures]
Theorem~\ref{thm:main} settles all four of Bala's conjectures at once, though only A and D
were still open. For B: $B(M)=0$ exactly when $M\le2$ (for $M\ge3$ the term $k=2$ or
$k=M-1$ contributes), so $v_{p}(a(n))>0$ iff $\lfloor n/p\rfloor\ge3$ iff $p\le n/3$. For
C: if $m=np+r$ with $0\le r<p$ then $\lfloor m/p^{i}\rfloor=\lfloor n/p^{i-1}\rfloor$ for
every $i\ge1$, independent of $r$, so $v_{p}(a(m))$ is constant across the block. For D:
taking $m=np$ in that same identity gives
$v_{p}(a(np))=B(n)+B(\lfloor n/p\rfloor)+B(\lfloor n/p^{2}\rfloor)+\cdots$. B and C are
credited on the entry to Adamczewski (2026); the derivations above are recorded only to
show that the closed form covers them.
\end{remark}

\begin{remark}[honest assessment]
The proof is short and uses nothing beyond Legendre's formula and two elementary floor
identities. It is an OEIS comment, not a paper. What makes it worth writing down is that
the closed form is not on the entry in any form, and that it replaces four separate
conjectures with one identity.

One caution the reader should carry: a benchmark of 492 Lean-formalised open OEIS
conjectures is currently being worked through by automated systems (Adamczewski,
arXiv:2608.11941, 12 August 2026, reporting 147 resolved; and earlier work on the same
492). Conjectures B and C of this entry fell to that effort three days before this note was
written. I have not been able to confirm whether A or D appear in that formalised set, and
anyone acting on this note should check the results repositories before assuming priority.
\end{remark}

\section{Verification}

A verification script, written separately from this note, checks every claim and prints
every number quoted. It (i) recomputes $a(1),\dots,a(27)$ from \eqref{eq:def} and matches
the published DATA of A129365 exactly; (ii) confirms $B(M)=0\iff M\le2$ for
$0\le M\le2000$, and prints $B(0),\dots,B(12)=0,0,0,1,1,4,3,8,8,12,13,22,17$;
(iii) compares $v_{p}(a(n))$ against $\sum_{i}B(\lfloor n/p^{i}\rfloor)$ for every
$1\le n\le60$ and every prime $p\le n$, that is 597 pairs $(n,p)$, with 0 mismatches; and
(iv) confirms directly that $a(n)$ is an integer for $1\le n\le100$.

\begin{thebibliography}{9}
\bibitem{oeis} N.~J.~A. Sloane et al., \emph{The On-Line Encyclopedia of Integer
Sequences}, \url{https://oeis.org/A129365}, entry A129365, consulted 23 August 2026; also
entries A092287 and A004125.
\bibitem{legendre} A.~M. Legendre, \emph{Th\'eorie des Nombres}, 3rd ed., Firmin Didot,
Paris, 1830.
\bibitem{adam} T.~Adamczewski, \emph{OEIS Open: How many conjectures can language models
turn into theorems?}, arXiv:2608.11941 [cs.AI], 2026.
\end{thebibliography}
\end{document}
